\documentclass[11pt]{article}
\usepackage[T1]{fontenc}
\usepackage{lmodern}
\usepackage[margin=1in]{geometry}
\usepackage{amsmath,amssymb,amsthm,mathtools,booktabs,array}
\usepackage[colorlinks=true,urlcolor=blue,citecolor=blue,linkcolor=blue]{hyperref}
\hypersetup{pdftitle={Switching reachability in twisted exclusion processes: extensions of Dabrowski, Poulain d'Andecy and Ragoucy's sector classification},pdfauthor={Haobo Ma; Wenlin Zhang}}
\newtheorem{theorem}{Theorem}
\newtheorem{lemma}[theorem]{Lemma}
\newtheorem{corollary}[theorem]{Corollary}
\theoremstyle{remark}
\newtheorem{example}[theorem]{Example}
\DeclareMathOperator{\Sym}{Sym}
\DeclareMathOperator{\supp}{supp}
\newcommand{\twist}{\mathcal{B}}
\newcommand{\charge}{\chi}
\title{Switching reachability in twisted exclusion processes:\\extensions of Dabrowski, Poulain d'Andecy\\and Ragoucy's sector classification}
\author{Haobo Ma\\\small ChronoAI Pte Ltd, Singapore\and Wenlin Zhang\\\small National University of Singapore\\\small ChronoAI Pte Ltd, Singapore}
\date{18 November 2026}
\begin{document}
\maketitle
\begin{abstract}
Extending Dabrowski, Poulain d'Andecy and Ragoucy's sector classification
\cite[Propositions~4.9, 4.10 and 4.14, pp.~13--15; arXiv:2602.18204v2,
11 June 2026]{DPR}, we answer their switching question
\cite[Section~7, p.~22]{DPR}: for a ring of at least three sites,
alternating two twisted symmetric exclusion processes reaches every
configuration precisely when the derived subgroup of the group generated by
their twists acts transitively on the local states.
For two sites, the criterion is instead transitivity of an explicitly generated
subgroup on ordered pairs. The proof combines independent coordinate
commutators with a conserved abelian charge. For three or more sites we classify
all components, without a transitivity assumption, by the orbit profile and a
charge in an explicit quotient of the twist group's abelianization.
The method also covers finite families of twists. Norris's positive-time
support theorem \cite[Theorem~3.2.1, p.~111]{Norris} gives the same
reachability for positive dwell times; periodic switching converges to the
uniform law on each component.
\end{abstract}
\noindent\small\textcopyright\ 2026 Haobo Ma and Wenlin Zhang.
Licensed under Creative Commons Attribution 4.0 International
(\href{https://creativecommons.org/licenses/by/4.0/}{CC BY 4.0}).\normalsize

\section{The question and the model}
For two twisted symmetric simple exclusion processes (SSEP), which permutations
allow repeated switches to give positive probability to every configuration?
This is the mathematical question in the Outlook of
Dabrowski, Poulain d'Andecy and Ragoucy \cite{DPR}. We give a necessary and
sufficient group-theoretic condition, including the distinct two-site case.
Their Proposition~5.2 supplies the branching probabilities; their
Sections~5.3--5.4 describe examples involving an identity twist or powers of a
common cycle. Their single-twist profile-and-charge classification
\cite[Proposition~4.9]{DPR} is the closest antecedent: we replace its cyclic
charge by an abelian quotient and use coordinate commutators to handle arbitrary
noncommuting twists. Their classification already appears in the first version,
released on 20 February 2026. The result here gives both full reachability and
all components.

Let $X=\{0,\ldots,N-1\}$, where $N\geq1$, and let $L\geq2$.
A configuration is $x=(x_1,\ldots,x_L)\in\Omega=X^L$.
For $s\in\Sym(X)$, the elementary moves exchange adjacent coordinates
$i,i+1$, $1\leq i<L$, and apply the boundary involution
\begin{equation}\label{eq:boundary}
 \twist_s(x_1,x_2,\ldots,x_{L-1},x_L)
 =(s^{-1}(x_L),x_2,\ldots,x_{L-1},s(x_1)).
\end{equation}
This is the direction of the boundary move in
\cite[Section~3.1, (3.1)--(3.3)]{DPR}.
The $s$-sectors are the connected components under these involutions.
They carry uniform stationary distributions
\cite[Proposition~4.4]{DPR}. For an $s$-sector $C$ and a $t$-sector $E$,
the probability of passing from the stationary distribution on $C$ to the
stationary distribution on $E$ after equilibration is
\begin{equation}\label{eq:branch}
 p(C,E)=\frac{|C\cap E|}{|C|}
 \qquad\text{\cite[Proposition~5.2]{DPR}}.
\end{equation}
In particular, this probability is positive exactly when $C\cap E\ne\varnothing$.

Fix twists $f,g$, and put
\[
 H=\langle f,g\rangle\leq\Sym(X),\qquad D=[H,H],\qquad
 K=\langle(f,f^{-1}),(g,g^{-1})\rangle\leq H\times H.
\]
Products act by composition, with the rightmost factor acting first;
$[u,v]=uvu^{-1}v^{-1}$. The action of $H\times H$ on $X^2$ is coordinatewise.
For $a\in X$, write $H_a=\{h\in H:h(a)=a\}$.
Let $\Gamma$ be the graph on $\Omega$ containing both sets of elementary moves.
Full switching reachability means that, from every sector of either model,
alternating equilibrated switches give positive probability to every
configuration after some finite number of switches.

\begin{theorem}\label{thm:main}
Full switching reachability is equivalent to connectivity of $\Gamma$.
For $L\geq3$, it is equivalent to transitivity of $D$ on $X$.
For $L=2$, it is equivalent to transitivity of $K$ on $X^2$, or equivalently to
\begin{equation}\label{eq:two}
 |K|=N^2\,|K\cap(H_a\times H_a)|
\end{equation}
for some, equivalently every, $a\in X$.
For any $L\geq2$, full switching reachability implies that $D$ is transitive.
\end{theorem}

\section{Switching support and coordinate operations}
For $s=f,g$ and $A\subseteq\Omega$, let $T_s(A)$ be the union of all $s$-sectors
meeting $A$. Equation~\eqref{eq:branch} shows that the support after each
equilibrated switch is obtained by applying the next $T_s$.
These operations are extensive: $A\subseteq T_s(A)$.
Consequently every elementary move of either type can be followed within the
next two switches: if its type is not next, first keep the current point in the
support, then follow the move at the following switch.
Thus every path in $\Gamma$ is eventually included in the alternating support.
Conversely, $T_f$ and $T_g$ preserve every component of $\Gamma$.
Starting with a nonempty sector in a component of size $m$, a simple path from
one of its points to any other point has at most $m-1$ edges. After at most
$2(m-1)$ switches, all points of the component occur simultaneously in the
support, since the supports are increasing. This proves the first equivalence
in Theorem~\ref{thm:main}, and gives a uniform bound.

Adjacent exchanges generate all site permutations. Composing the boundary
move with the exchange of the first and last sites, and conjugating by site
permutations, therefore gives, for every $i\ne j$, the coordinate operation
\[
 E_{ij}(s):\quad x_i\longmapsto s(x_i),\quad
 x_j\longmapsto s^{-1}(x_j),\quad x_k\longmapsto x_k\ (k\ne i,j),
 \qquad s=f,g.
\]
Let $R\leq H^L$ be the subgroup generated by these operations.
Every element of $R$ is realized by a path in $\Gamma$ from each configuration.
For fixed $i\ne j$, the subgroup generated by $E_{ij}(f),E_{ij}(g)$ is the
copy of $K$ on coordinates $i,j$; its projection to coordinate $i$ is $H$.

\begin{lemma}\label{lem:coordinate}
If $L\geq3$, then $D^L\leq R$.
\end{lemma}
\begin{proof}
Fix $i$ and choose distinct $j,k$ different from $i$.
Given $u,v\in H$, surjectivity of the pair projections provides an element of
$R$ supported on $i,j$ whose $i$-coordinate is $u$, and an element supported
on $i,k$ whose $i$-coordinate is $v$.
Their commutator is $[u,v]$ at $i$ and the identity at every other coordinate:
at $j$ or $k$, one of the two factors is the identity.
Since $D$ is generated by all such commutators, $R$ contains $D$ at coordinate
$i$ alone. Multiplying these subgroups for all $i$ gives $D^L\leq R$.
\end{proof}
If $D$ is transitive, choose independently $d_i\in D$ carrying $x_i$ to $y_i$.
Lemma~\ref{lem:coordinate} realizes $(d_1,\ldots,d_L)$, so $\Gamma$ is connected.
This proves sufficiency for $L\geq3$.

\section{The charge obstruction and all components}
An elementary move preserves the number of coordinates in each $H$-orbit on
$X$. Hence connectivity forces $H$ to be transitive: otherwise the constant
configurations at points in different $H$-orbits cannot be connected.

Suppose now that $H$ is transitive. Fix $a\in X$ and put
\[
 M=D H_a,\qquad A=H/M.
\]
The subgroup $M$ is normal, because its image in the abelian group $H/D$ is a
subgroup. Thus $A$ is a finite abelian group, written additively.
For $b=h(a)$ define
\begin{equation}\label{eq:charge}
 \charge(b)=hM\in A,\qquad Q(x)=\sum_{i=1}^{L}\charge(x_i).
\end{equation}
This is well defined: two choices of $h$ differ on the right by an element of
$H_a\leq M$. Moreover
\begin{equation}\label{eq:equivariance}
 \charge(u(b))=uM+\charge(b)\qquad(u\in H).
\end{equation}
The fibers of $\charge$ are exactly the $D$-orbits. Indeed, equal cosets give
$h_2^{-1}h_1\in D H_a$, so $h_1(a)$ and $h_2(a)$ are in the same $D$-orbit by
normality of $D$; the converse follows by the same calculation.
In particular, $|A|$ is the number of $D$-orbits, and $A$ is trivial precisely
when $D$ is transitive.

Site exchanges preserve $Q$. Under~\eqref{eq:boundary}, the charges at the two
ends become $-sM+\charge(x_L)$ and $sM+\charge(x_1)$, whose sum is unchanged.
Thus $Q$ is constant on every component of $\Gamma$ for every $L\geq2$.
If $A$ is nontrivial, surjectivity of $\charge$ gives $b\in X$ with
$\charge(b)\ne0=\charge(a)$, so $(a,\ldots,a)$ and $(b,a,\ldots,a)$ have
different charges. This proves the necessity assertion in
Theorem~\ref{thm:main}, including for $L=2$.

\begin{theorem}\label{thm:components}
Suppose $H$ is transitive and $L\geq3$. Two configurations lie in the same
component of $\Gamma$ if and only if their total charges $Q$ agree.
Consequently $\Gamma$ has exactly $q=|H/(D H_a)|$ components, each of size
$N^L/q$.
\end{theorem}
\begin{proof}
Necessity was proved above. For sufficiency, take $Q(x)=Q(y)$.
For any $(u,v)\in K$, the sum $uM+vM$ is zero: the map
$H\times H\to A$, $(u,v)\mapsto uM+vM$, is a homomorphism and vanishes on
both generators of $K$.
Its first projection being surjective, $K$ therefore realizes, on any distinct
coordinates $i,L$, any prescribed charge increment $c\in A$ at $i$ and $-c$
at $L$. Apply such operations successively for $i=1,\ldots,L-1$ to make the
first $L-1$ charges agree with those of $y$.
Conservation of $Q$ then makes the last charges agree as well.
The resulting entries lie in the same respective $D$-orbits as the entries of
$y$. Lemma~\ref{lem:coordinate} changes them independently to the entries of
$y$. This proves the classification.

Because $D$ is normal and $H$ is transitive, all its orbits have the same size,
namely $N/q$. Each element of $A$ occurs as a total charge: choose the first
$L-1$ individual charges arbitrarily and let the last enforce the required
sum. There are $q^{L-1}$ such charge tuples, each representing $(N/q)^L$
configurations. Every component therefore has size $N^L/q$.
\end{proof}

The next statement removes the transitivity assumption. It is the full
component classification; Theorem~\ref{thm:components} is its one-orbit case.

\begin{theorem}\label{thm:profile}
Let $L\geq3$, and let $O_1,\ldots,O_r$ be the $H$-orbits on $X$.
Choose $a_j\in O_j$ and fix a profile $n=(n_1,\ldots,n_r)$ with
$n_j\geq0$ and $\sum_j n_j=L$. Put $B=H/D$, let $B_j$ be the image
of $H_{a_j}$ in $B$, and define the finite abelian group
\[
 C_n=B\Big/\sum_{j:n_j>0}B_j.
\]
For $b=h(a_j)\in O_j$ with $n_j>0$, set $\chi_n(b)$ equal to the image of $h$ in $C_n$.
Within the configurations having $n_j$ entries in $O_j$, the connected
components are precisely the fibers of $Q_n(x)=\sum_i\chi_n(x_i)$.
There are $q_n=|C_n|$ such components, each of cardinality
\begin{equation}\label{eq:profile-size}
 \frac{1}{q_n}\frac{L!}{\prod_j n_j!}\prod_j |O_j|^{n_j}.
\end{equation}
\end{theorem}
\begin{proof}
The stabilizers of occupied orbits have zero image in $C_n$, so $\chi_n$ is well defined.
The profile and $Q_n$ are preserved by every move, by the same cancellation
as in~\eqref{eq:equivariance}.
For sufficiency, first permute sites so two configurations with the same
profile have the same orbit at each site. Write $j(i)$ for that orbit's index.
The $D$-orbits within $O_j$ are identified with $B/B_j$:
if $hD$ has image in $B_j$, then $h\in D H_{a_j}$, and normality of $D$
gives exactly the $D$-orbit relation.
Let $b_i\in B$ lift the desired change in $B/B_{j(i)}$ at site $i$.
Equality of $Q_n$ says that $\sum_i b_i$ belongs to
$\sum_i B_{j(i)}$. Subtract appropriate elements of $B_{j(i)}$ from the
individual lifts to arrange $\sum_i b_i=0$ without changing those desired
changes. The pair operations $E_{iL}(f),E_{iL}(g)$ generate every increment
$(b,-b)$ in $B$ at sites $i,L$: their first increments generate $B$ and
their second increments are the negatives. Apply these operations for
$i<L$ to realize $(b_1,\ldots,b_L)$. Each site now lies in the desired
$D$-orbit. Lemma~\ref{lem:coordinate} finishes independently at each site.

For fixed site orbits, the map from local states to $\prod_i B/B_{j(i)}$
has constant fiber size: normality of $D$ makes its orbits within each
$O_j$ equicardinal. The sum map from this product onto $C_n$ is a
surjective homomorphism, so all its fibers have equal size. There are
$L!/\prod_j n_j!$ arrangements of the site orbits. Dividing their total
number of configurations by $q_n$ gives~\eqref{eq:profile-size}.
\end{proof}

\begin{corollary}\label{cor:cyclic}
If $f,g$ are powers of a permutation $h$, write $H=\langle f,g\rangle$
and let its orbits have lengths $c_1,\ldots,c_r$.
For each profile $n$, put $d_n=\gcd(c_j:n_j>0)$.
For $L\geq3$ there are exactly $d_n$ components with that profile,
each of size $\frac{L!}{\prod_j n_j!}\prod_j c_j^{n_j}/d_n$.
In particular, when $H$ is transitive, there are $N$ components of size
$N^{L-1}$.
\end{corollary}
\begin{proof}
The cyclic group $H$ has order $m=\operatorname{lcm}(c_1,\ldots,c_r)$.
Its stabilizer on an orbit of length $c_j$ is generated by the $c_j$th
power of a generator of $H$. The sum of these stabilizers therefore has
index $d_n$. Apply Theorem~\ref{thm:profile}.
\end{proof}
These formulas reprove the profile-and-charge sector counts and sizes of
\cite[Propositions~4.9, 4.10 and 4.14]{DPR} for the equivalent cyclic model
in this range, rather than asserting them as new. For a single twist those
authors also cover $L=2$.

\section{Two sites and consequences}
For $L=2$ let $J$ exchange the coordinates. The elementary boundary maps are
$\twist_s=E_{21}(s)J$. Exchanging the two factors sends each generator of $K$
to its inverse, so $J$ normalizes $K$. The full move group is consequently
generated by $K$ and $J$.
Since $J$ fixes $(a,a)$, the orbit of this point under the full move group is
exactly its $K$-orbit. Thus connectivity forces this $K$-orbit to be $X^2$;
conversely, transitivity of $K$ clearly suffices for connectivity.
The stabilizer of $(a,a)$ in $K$ is $K\cap(H_a\times H_a)$.
The bijection between its left cosets and the orbit of $(a,a)$ gives the
orbit size $|K|/|K\cap(H_a\times H_a)|$.
This is $N^2$ exactly in~\eqref{eq:two}, completing the proof of
Theorem~\ref{thm:main}.

\begin{example}\label{ex:three}
Take $X=\{0,1,2\}$, $f=(0\,1)$ and $g=(0\,1\,2)$.
Then $H=\Sym(X)$ and $D=\langle g\rangle$ is transitive, so every $L\geq3$
has full switching reachability. For two sites, however,
\[
 K=\{(h,fhf^{-1}):h\in H\},\qquad |K|=6<9.
\]
Indeed, conjugation by $f$ fixes $f$ and inverts $g$, so the displayed subgroup
has precisely the required generators. Its action preserves the condition
$y=f(x)$, splitting $X^2$ into components of sizes $3$ and $6$:
the size-$3$ component is $\{(0,1),(1,0),(2,2)\}$.
The complement is the other component, since the diagonal action of $S_3$
is transitive on unequal ordered pairs after replacing $y$ by $f^{-1}(y)$.
Thus the restriction $L\geq3$ in the derived-group sufficiency statement is sharp.
\end{example}

\begin{corollary}\label{cor:commuting}
If $N>1$ and $f,g$ commute, full switching reachability fails for every $L\geq2$.
In particular, this holds if one twist is the identity or both are powers of
the same permutation.
\end{corollary}
\begin{proof}
The group $H$ is abelian, so $D$ is trivial and is not transitive on $X$.
Apply the necessity part of Theorem~\ref{thm:main}.
\end{proof}
This distinguishes the spreading and splitting examples in
\cite[Sections~5.3--5.4]{DPR} from unrestricted reachability.
For $N=1$ there is a single configuration, and all the criteria hold,
including~\eqref{eq:two}.

\begin{corollary}\label{cor:finite}
Replace equilibration by evolution for any strictly positive time at each
alternating stage. The switching support, the criteria of
Theorem~\ref{thm:main}, and the bound $2(m-1)$ on the number of switches needed
to fill an initially occupied component of size $m$ remain valid.
\end{corollary}
\begin{proof}
Let $U_1,\ldots,U_L$ be the permutation matrices for the $L-1$ adjacent
exchanges and the boundary involution for a fixed twist $s$.
Its Markov generator is $M_s=\sum_{r=1}^{L}U_r-LI$.
For $t>0$,
\begin{equation}\label{eq:semigroup}
 e^{tM_s}=e^{-Lt}\sum_{n=0}^{\infty}\frac{t^n}{n!}
             \left(\sum_{r=1}^{L}U_r\right)^n.
\end{equation}
Every summand has nonnegative entries, and a given entry is positive in some
power precisely when an elementary-move path joins its two configurations;
the $n=0$ term includes each configuration itself. Hence $e^{tM_s}$ has
positive entries exactly within $s$-sectors. Its action on the support is
therefore exactly $T_s$. The alternating-support argument applies unchanged.
Equation~\eqref{eq:semigroup} is the finite-state exponential-series
construction of a continuous-time chain; see also \cite[Section~20.1]{LP}.
The stagewise support assertion is a specialization of the earlier general
support theorem of Norris \cite[Theorem~3.2.1, p.~111]{Norris}.
\end{proof}
This establishes qualitative finite-time accessibility. Quantitative transition
probabilities, relaxation rates, and randomly varying twists involve further
questions, as do general set-theoretical Yang--Baxter models with different moves.

\begin{corollary}\label{cor:family}
Let $S$ be a nonempty finite family of twists and let $H=\langle S\rangle$.
For $L\geq3$, full switching reachability is equivalent to transitivity of
$[H,H]$, and Theorem~\ref{thm:profile} classifies all components.
These conclusions hold with either equilibration or strictly positive dwell
times whenever every twist occurs infinitely often. If every consecutive
block of $b$ stages contains every twist, an initial nonempty support fills
its component of size $m$ within at most $b(m-1)$ stages.
\end{corollary}
\begin{proof}
Replace the two pair generators by $E_{ij}(s)$ for $s\in S$.
Their projections generate $H$, so Lemma~\ref{lem:coordinate} and the
charge proofs apply verbatim. Each support operation is extensive, and
every edge type eventually occurs. Thus a finite path can be followed in
the prescribed order; under the block assumption each edge costs at most
$b$ stages. A simple path has at most $m-1$ edges. Positive-time support
is the same by~\eqref{eq:semigroup}.
\end{proof}

\begin{corollary}\label{cor:limit}
Repeat a fixed nonempty finite list of twists, using either equilibration at each stage
or a fixed strictly positive dwell time for each stage of the period.
For any $L\geq2$, the distribution converges to the mixture of uniform laws
on the union-graph components, with weights equal to their initial masses.
\end{corollary}
\begin{proof}
Each stage matrix is doubly stochastic, has strictly positive diagonal,
and has positive entries precisely within its sectors. Restrict to a
union component $C$ of size $m$ and let $P$ be one period's transition matrix.
Every edge of a simple path in $C$ can be traversed during one period,
waiting at all other stages. Hence $P^k$ is strictly positive on $C$ for
$k=\max(1,m-1)$. Let $\epsilon$ be its smallest entry there and let $U$
have every entry $1/m$. If $m\epsilon<1$, then
$P^k=m\epsilon U+(1-m\epsilon)V$ for a doubly stochastic matrix $V$.
Thus every block of $k$ periods contracts the distance in total variation
to the uniform law by at most $1-m\epsilon<1$. If $m\epsilon=1$,
$P^k=U$ directly. Intermediate stages preserve the uniform law and do not
increase that distance. Component masses are invariant, giving the stated
mixture.
\end{proof}

\section{Source coverage and verification}
The following table distinguishes the question answered here from the source's
other directions.
\begin{center}\small
\begin{tabular}{@{}p{0.33\textwidth}p{0.62\textwidth}@{}}\toprule
Source item & Coverage \\\midrule
\cite{DPR}, Section~7: switching reachability &
Answered: yes; Theorem~\ref{thm:main} gives necessary and sufficient
conditions, including the two-site case.\\[3pt]
\cite{DPR}, Propositions~4.9, 4.10, 4.14: cyclic sectors &
Reproved; prior proof Dabrowski, Poulain d'Andecy and Ragoucy (2026).
Corollary~\ref{cor:cyclic} recovers the formulas for $L\geq3$;
Theorem~\ref{thm:profile} extends them to noncyclic twist groups.\\[3pt]
\cite{DPR}, Section~5.2: spreading, splitting and oscillations &
Prior formulas cited in~\eqref{eq:branch}; support characterized above.
Corollary~\ref{cor:limit} proves the limiting law under repeated fixed periods
(not posed in this precise form).\\[3pt]
\cite{DPR}, Sections~5.3--5.4: identity and power twists &
Corollary~\ref{cor:commuting} excludes full reachability for $N>1$.\\[3pt]
\cite{DPR}, Section~7: finite-time effects &
Stage support reproved; prior proof Norris (1997), Theorem~3.2.1.
Corollaries~\ref{cor:finite} and~\ref{cor:family} give the switching
consequences (not posed in this precise form); quantitative and
random-time dynamics not covered.\\[3pt]
\cite{DPR}, Remark~6.3: the three-species non-equivalence &
Prior result cited: Dabrowski, Poulain d'Andecy and Ragoucy
\cite[Remark~A.8: v1, p.~30, 21 July 2026; v2, p.~31]{DPRreactive};
not a consequence of the switching criterion.\\[3pt]
\cite{DPR}, Section~7: integrable boundaries and deformations &
Not covered by this criterion. Reactive models and boundaries are studied in
\cite{DPRreactive}.\\[3pt]
\cite{DPR}, Remark~3.5: permutation roots; Remarks~5.6--5.7: identities &
Prior results cited there; no new root criterion or multinomial identity
is asserted here.\\\bottomrule
\end{tabular}
\end{center}

The pinned Lean source is
\href{https://github.com/the-omega-institute/trureturing/blob/6bfc394bbd5147de839dd2116b985b438fcae8cf/D5/S3/StatisticalMechanics/ParticleProcesses/SwitchedTwistedSSEP.lean}{\texttt{SwitchedTwistedSSEP.lean}}.
Its declaration \texttt{result} states the all-parameter switching classification,
the $L\geq3$ derived-group criterion, the $L=2$ stabilizer criterion,
positivity of the sector-overlap probabilities, and equality of the two
alternating-support definitions. These statements correspond to
Theorem~\ref{thm:main} and the equilibrated-support argument.
The profile-and-charge classification in Theorem~\ref{thm:profile}, its
transitive specialization in Theorem~\ref{thm:components}, the explicit
example, and Corollaries~\ref{cor:cyclic}, \ref{cor:finite},
\ref{cor:family} and \ref{cor:limit} have written proofs here. They are not additional claims
of the declaration \texttt{result}. The accompanying exact-arithmetic program independently
checks finite instances of the group criteria, component counts and sizes,
complete profile-and-charge fibers, branching probabilities, and alternating supports;
its range and results are
given in \texttt{verification.json}. The all-parameter proofs do not depend
on that finite range.

\begin{thebibliography}{9}
\bibitem{DPR}
M.~Dabrowski, L.~Poulain d'Andecy and E.~Ragoucy,
\emph{Twisted symmetric exclusion processes and set-theoretical $R$-matrices},
J.~Phys.~A: Math.~Theor. \textbf{59} (2026), 245203.
\href{https://doi.org/10.1088/1751-8121/ae751a}{doi:10.1088/1751-8121/ae751a}.
\href{https://arxiv.org/abs/2602.18204v2}{arXiv:2602.18204v2} (11 June 2026).
First version: \href{https://arxiv.org/abs/2602.18204v1}{arXiv:2602.18204v1}
(20 February 2026).

\bibitem{DPRreactive}
M.~Dabrowski, L.~Poulain d'Andecy and E.~Ragoucy,
\emph{Integrable multi-species SSEP with reactive particle species},
J.~Stat.~Mech. (2026), 113102.
\href{https://doi.org/10.1088/1742-5468/aeb496}{doi:10.1088/1742-5468/aeb496}.
\href{https://arxiv.org/abs/2607.18959v2}{arXiv:2607.18959v2} (12 October 2026).
First version: \href{https://arxiv.org/abs/2607.18959v1}{arXiv:2607.18959v1}
(21 July 2026).

\bibitem{LP}
D.~A.~Levin and Y.~Peres, with contributions by E.~L.~Wilmer,
\emph{Markov Chains and Mixing Times}, second edition,
American Mathematical Society, 2017.
\href{https://doi.org/10.1090/mbk/107}{doi:10.1090/mbk/107}.

\bibitem{Norris}
J.~R.~Norris, \emph{Markov Chains}, Cambridge University Press, 1997.
\href{https://doi.org/10.1017/CBO9780511810633}{doi:10.1017/CBO9780511810633}.
Section~3.2, Theorem~3.2.1, p.~111.
\end{thebibliography}
\end{document}
